\section{问题三的模型建立与求解}
\label{sec:question-three}

\subsection{问题分析与模型输入}
问题二确定了每张图在各核数下的分核与核内顺序；问题三增加共享只读Cache，关键变化是同一份张量在不同核心上的读取时刻和读取来源。为分开观察硬件变化与调度调整，本章以同图同核数下问题二的固定核数方案$X_B^*$为起点，包括操作到子图的映射、子图到核心的映射和各核内顺序。先在C配置下评价这份计划，再围绕共享输入产生新候选，由此得到$F_B(X_B^*)$、$F_C(X_B^*)$和最终选择的$F_C(X_C^*)$。500组配对记录中有497组的C初始计划哈希与$X_B^*$一致；case\_054的5核、case\_062的4核和case\_085的5核保留了实验记录中的初始计划差异，后文将其单独标注。

共享只读Cache容量为1048576 B，读取带宽为250 B/cycle；DDR带宽仍为60 B/cycle，两个读取池独立。Cache初始为空，按先进先出淘汰。逻辑张量标识用于命中查询，实际COPY事件来自问题二的完整核内展开，Spill产生的物理版本也由评估程序处理。每组配对都记录三次评价的计划哈希，并以哈希一致的497组作为严格同计划解释的依据。
\begin{table}[H]\centering\caption{问题三的固定输入和配对量}\label{tab:c-inputs}
\begin{tabularx}{\linewidth}{lY}\toprule
对象 & 值或判定条件\\\midrule
同图同核起点 & 问题二固定核数方案$X_B^*$；497/500组C初始计划哈希一致\\
共享Cache & 容量1048576 B；读带宽250 B/cycle；初始队列为空\\
外部存储 & DDR带宽60 B/cycle；写出与未命中读仍走DDR\\
正式响应 & 总工期、结构与恢复搬运、命中和未命中字节、完整计划\\
选择预算 & 常规额度为$\max\{\min(q,2),m+1\}$；$m$为去重后必评锚点数\\\bottomrule
\end{tabularx}\end{table}

\subsection{基于查询与完成事件的FIFO模型}
设$\mathcal C(\tau)$为时刻$\tau$按插入次序排列的逻辑张量队列。合格读入事件$e$在发射时刻$q_e$查询$t_e$，其读取来源由
\begin{equation}
H_e=\mathbf 1\{t_e\in\mathcal C(q_e)\},\qquad
B_e=\begin{cases}250,&H_e=1,\\60,&H_e=0\end{cases}\ \mathrm{B/cycle}
\label{eq:c-query}
\end{equation}
确定。命中读进入Cache池，未命中读进入DDR池；命中只改变本次读取所用的服务池，不把条目移到队尾，查询也不改变FIFO次序。在该COPY\_IN完成时，评估器再尝试把张量插入队尾：若当前仍在队列中则不改变次序；若已经被其他完成事件淘汰，则可能重新插入；若大小超过容量则不插入；空间不足时先逐项淘汰队首。同一时刻的完成、插入和新发射按评估程序的固定事件次序处理，尚未完成的在途读不会被后续查询命中。式\eqref{eq:c-query}中的$B_e$是该次读取所用的服务带宽。

考虑容量为$2b$、三个张量$x,y,z$大小均为$b$的符号例。已有队列$[x,y]$，对$x$的命中不会将其移到队尾。随后$z$完成并插入，队列变为$[y,z]$；如果先前对$x$的在途读取此时才完成，它又能在队尾重新插入，队列变为$[z,x]$。同一组张量的访问次数并不足以确定这段队列演化，完成事件的先后与查询先后都必须考虑。

更关键的是并发查询。两个核心在首个读入完成前向空Cache发射相同张量的请求，两次都未命中；只有后发请求在首次完成插入之后查询且条目未被淘汰，才可能命中。记合格读入总命中字节为$Q_L$，未命中字节为$Q_D$，则报告的字节命中率为
\begin{equation}
h_{\rm byte}=\frac{Q_L}{Q_L+Q_D}.
\label{eq:hit-rate}
\end{equation}
分母为零时取零。该指标只描述读取来源，额外搬运仍按实际生成的COPY事件统计。

DDR和Cache两个共享服务池各自按活动请求分配参考工作量。若池$p$中同时有$m_p(\tau)$个未完成请求，请求$e$的剩余参考工作量$R_e$满足
\begin{equation}
\frac{\mathrm dR_e}{\mathrm d\tau}=-\frac{1}{m_p(\tau)},\qquad
R_e(q_e)=\max\left\{1,\left\lceil\frac{b_e}{B_e}\right\rceil\right\}.
\label{eq:c-service}
\end{equation}
一次命中既可能减轻DDR池竞争，也可能改变后继操作的发射时刻，从而改变关键路径、另一池负载及后续FIFO状态。因此即使$h_{\rm byte}>0$，总工期仍可能不变。

\begin{algorithm}[H]\caption{评估程序中只读FIFO Cache的事件规则}\label{alg:c-fifo}
\begin{algorithmic}[1]
\REQUIRE 已展开的核级执行图、COPY事件、Cache容量及两类带宽
\ENSURE 完成时间线、Cache事件、命中与未命中字节
\STATE 初始化空FIFO队列、DDR池与Cache读池
\WHILE{仍有操作未完成}
\STATE 按评估程序的事件次序处理当前完成事件，推进相应池的剩余工作量
\FOR{当前完成的每个合格COPY\_IN}
\IF{逻辑张量不在队列且大小不超过容量}
\STATE 从队首淘汰至空间充足，并插入该张量
\ENDIF
\ENDFOR
\STATE 更新依赖、流水线与跨核释放就绪状态
\FOR{当前可发射的合格COPY\_IN}
\STATE 在发射时查询队列，选择Cache读池或DDR池并累计字节
\ENDFOR
\STATE 重排两池的预计完成事件并推进事件时间
\ENDWHILE
\end{algorithmic}\end{algorithm}
\subsection{Cache查询时序的解析算例}
为把式\eqref{eq:c-query}与事件先后联系起来，构造一个不参与100图统计的解析小例：三个核心各需读同一份60000 B张量$x$，Cache容量足以容纳$x$。独占DDR读取为$60000/60=1000$ cycle，独占Cache读取为$60000/250=240$ cycle。若三核在时刻0同时发射，三次查询都看到空队列，三个DDR请求平分服务，读取均在时刻3000完成，命中字节为0。若第一核在时刻0发射并于1000完成插入，其余两核在时刻1000才发射且$x$未被淘汰，则两次都命中；两个Cache请求平分读池，各需480 cycle，读取于1480完成。表\ref{tab:c-query-example}把发射、插入和完成三个状态明确分开。
\begin{table}[H]\centering\caption{三核重复读取60000 B张量的解析时序}\label{tab:c-query-example}
\begin{tabularx}{\linewidth}{lYrrr}\toprule
情形 & 发射与插入顺序 & DDR读取 & Cache读取 & 最后一次读完成/cycle\\\midrule
同时查询 & 三核均在0发射，首次插入发生在完成时 & 3次 & 0次 & 3000\\
错峰查询 & 第一核0发射、1000插入，其余两核1000发射 & 1次 & 2次 & 1480\\\bottomrule
\end{tabularx}\end{table}
\begin{figure}[H]\centering
\begin{tikzpicture}[x=1.65cm,y=.9cm,>=Stealth,every node/.style={font=\small}]
\draw[->] (0,0)--(3.5,0) node[right] {$\tau/1000$ cycle};
\foreach \x/\lab in {0/0,1/1,1.48/1.48,3/3}{\draw (\x,.06)--(\x,-.06) node[below] {\lab};}
\node[left] at (0,1.0) {错峰查询};
\node[left] at (0,2.4) {同时查询};
\draw[blue!65!black,very thick] (0,2.4)--(3,2.4);
\node[above] at (1.5,2.4) {3次DDR未命中并发};
\draw[blue!65!black,very thick] (0,1.0)--(1,1.0);
\node[above] at (.5,1.0) {1次DDR};
\draw[green!55!black,very thick] (1,1.0)--(1.48,1.0);
\node[anchor=west] at (1.58,1.0) {2次Cache命中};
\draw[dashed] (1,.12)--(1,2.9);
\node[anchor=west] at (1.65,3.25) {首读完成并插入：$\tau=1000$};
\end{tikzpicture}
\caption{相同访问次数下的两种查询顺序；仅为解析示例}\label{fig:c-query-timeline}
\end{figure}
错峰发射在真实计划中须由前置计算、依赖或队列中的其他操作形成，这些操作本身也占用资源。

\subsection{共享输入方案与求解}
\textbf{Step 1：生成候选。}以$X_B^*$为种子，尝试迁移、同核合并、核内调序以及跨核分散，并查找跨核重复读取的逻辑张量。对按张量大小和潜在重复读量排序的有限张量集合，枚举相关消费者组或节点在现有核心序列中的前移、后移和跨核插入位置；对满足容量条件的共享输入，再指定一个领读核，将其消费者移至该核序列前部，并把其他消费者后移到各自序列尾部，形成错峰查询候选。候选只保留合法映射与顺序；共享张量集合随当前Cache容量过滤，因而容量扰动会同步改变候选组合。不合法的构造不进入完整评价。

\textbf{Step 2：轻量排序。}按路径递推、最大核心计算量、结构COPY的DDR参考工作量和容量压力估计候选代价；C另外在收缩后的商图上回放逻辑COPY查询的相对时序，记录查询时的FIFO命中代理、未命中量和潜在Cache服务工作量。该回放保留每核顺序、跨核释放和完成后插入规则，但不模拟服务池的并发；正式工期、命中和淘汰由问题三的事件模拟给出。

\textbf{Step 3：完整评价与选优。}把B最终计划$X_B^*$、整图、分量装桶和已评价的低核数方案中去重后的方案置于必评锚点，其余C候选按轻量分数排列。设去重后必评数为$m$、请求额度为$q$，常规完整调用额度为$\max\{\min(q,2),m+1\}$。展开失败只记录错误，不占完整调用额度；完整评价失败占用一次额度。首个成功结果记为初始方案，最终按工期、额外搬运的字典序在全部成功结果中选定$X_C^*$。得到首个C结果后，从其Cache事件中提取重复未命中张量，生成事件锚定的合法调序候选；若仍有额度，再按候选代表补评。下文500组B/C配对按记录中的B计划与C初始计划计算三段工期；其中497组哈希一致，用于严格同计划效应分解，与各核数的独立最终选择曲线分开报告。

\begin{algorithm}[H]\caption{问题三候选选择与同计划对照}\label{alg:c-select}
\begin{algorithmic}[1]
\REQUIRE 问题二固定核数方案$X_B^*$、图、核数、固定硬件参数及请求完整调用额度$q$
\ENSURE 问题三选中计划$X_C^*$及三段工期比
\STATE 保留$X_B^*$，生成迁移、同核合并、调序、分散及领读核错峰代表
\STATE 逐个检查合法性、按完整计划去重，计算轻量排序分数
\STATE 将$X_B^*$、整图、分量装桶和已评价的低核数方案置于去重后的必评队列
\STATE 令$m$为必评数，完整调用额度为$\max\{\min(q,2),m+1\}$
\WHILE{未达到常规完整调用额度且还有未尝试候选}
\STATE 取必评队列或轻量排序最靠前的未尝试候选，先展开检查再做C场景事件模拟
\STATE 展开失败不占额度；完整评价保存成功结果或完整错误记录
\ENDWHILE
\STATE 从首个成功结果提取重复未命中事件，生成事件锚定调序候选并在剩余额度内补评
\STATE 在成功结果中按工期、额外搬运的字典序选$X_C^*$
\STATE 核对C初始计划与$X_B^*$一致，计算式\eqref{eq:cache-factors}
\end{algorithmic}\end{algorithm}

\subsection{共享输入聚合与错峰的解析边界}
仍用60000 B共享输入，令五个V流水线操作各计算1000 cycle，暂不计输出。下面比较两种给定的局部请求时序下服务池的边界工期。全部聚到一核，只读一次DDR却顺序完成五个计算，理论完成时间为$1000+5\times1000=6000$ cycle。五核分散且同时查询空Cache时，五次未命中读平分DDR，读阶段持续$5\times60000/60=5000$ cycle，再各算1000 cycle，总时间同为6000。若只有第一核在0查询，其余四核由外部前置关系约束到1000才查询，四次命中合计240000 B，Cache读阶段需$240000/250=960$ cycle，后四核在2960完成。读取时序与Cache服务能力共同决定局部读计算完成时间。
\begin{table}[H]\centering\caption{五个V操作共享60000 B输入的解析边界（非正式实验）}\label{tab:c-five-example}
\begin{tabularx}{\linewidth}{lYr}\toprule
结构 & 读取与执行条件 & 解析工期/cycle\\\midrule
单核聚合 & 一次DDR读取，五次V操作串行 & 6000\\
五核分散、无Cache & 五次DDR读同时开始、共享60 B/cycle & 6000\\
五核分散、Cache同时查询 & 空队列下五次均未命中 & 6000\\
五核分散、其余四核错峰 & 后四次命中，Cache带宽250 B/cycle & 2960\\
同上、带宽减半 & 后四次命中，Cache带宽125 B/cycle & 3920\\
同上、带宽加倍 & 后四次命中，Cache带宽500 B/cycle & 2480\\
同上、容量不足60000 B & 张量不能插入，后四次仍未命中 & 6000\\\bottomrule
\end{tabularx}\end{table}
五核分散在250 B/cycle条件下，第一核于2000完成计算；其余四核从1000起共读960 cycle，再计算1000 cycle，于2960完成。带宽减半和加倍时，将960分别改为1920、480，即得3920、2480。

\subsection{同计划对照与调度适配}
同一图同一核数取得三个工期：$F_B(X_B^*)$、$F_C(X_B^*)$、$F_C(X_C^*)$。定义
\begin{equation}
\begin{split}
S_{\rm hw}&=\frac{F_B(X_B^*)}{F_C(X_B^*)},\\
S_{\rm adapt}&=\frac{F_C(X_B^*)}{F_C(X_C^*)},\\
S_{\rm opt}&=\frac{F_B(X_B^*)}{F_C(X_C^*)}=S_{\rm hw}S_{\rm adapt}.
\end{split}
\label{eq:cache-factors}
\end{equation}
第一个比值改变Cache配置而保持计划，第二个比值在C配置下比较初始计划与最终选择，最后一个比值描述从B计划到C最终计划的综合变化。500条配对记录的乘积恒等式最大数值残差为$2.22\times10^{-16}$；其中497条记录的B计划哈希与C初始计划哈希一致，3条记录的差异已在本节开头列出。因此，表\ref{tab:cache-factors}的均值按全部500条原始配对记录计算，严格同计划的局部解释采用哈希一致子集。
\begin{table}[H]\centering\caption{每核100图的Cache收益三段分解（raw final\_selection）}\label{tab:cache-factors}
\begin{tabular}{crrrrr}\toprule
核数 & 硬件效应均值 & 适配效应均值 & 综合效应均值 & 硬件改善图数 & 适配改善图数\\\midrule
1 & 1.0023 & 1.0145 & 1.0170 & 27 & 49\\
2 & 1.0079 & 1.0142 & 1.0227 & 31 & 47\\
3 & 1.0090 & 1.0127 & 1.0220 & 40 & 46\\
4 & 1.0164 & 1.0092 & 1.0257 & 40 & 32\\
5 & 1.0146 & 1.0139 & 1.0289 & 48 & 38\\\bottomrule
\end{tabular}\end{table}
图\ref{fig:cache-direct-comparison}给出$F_B(X_B^*)$、$F_C(X_B^*)$和$F_C(X_C^*)$的绝对平均工期及后二者相对$F_B(X_B^*)$的逐图比值均值；绝对工期受图规模影响，比值均值才反映硬件效应与调度适配的相对幅度。
\samplefigure{cache_direct_comparison}{问题三无Cache、固定计划Cache与重调度Cache的同核直接比较}{fig:cache-direct-comparison}

图\ref{fig:cache-coverage-radar}给出每核100图中工期严格缩短的占比。五核的硬件效应、C内适配与综合效应分别在48、38和65条原始配对记录中大于1；在哈希一致的98个五核配对中，三类记录分别为48、37和64。硬件效应与适配效应的计数存在交集，说明Cache读服务和计划调整作用于不同的关键路径位置。
\samplefigure{cache_effect_coverage_radar}{五个核数下硬件、适配与综合工期改善的逐图覆盖率分组条形图}{fig:cache-coverage-radar}

五核$C(X_B^*)$的字节加权命中率为22.65\%，500条配对中326条最终记录包含命中、174条命中字节为零。命中并不自动转化为尾部工期缩短：在哈希一致的五核配对中，29条记录虽有命中但$S_{\rm hw}=1$，另有34条记录同时保持零命中和零硬件效应。
\samplefigure{cache_hw_effect_pie}{五核硬件效应的图数构成及改善幅度分类条形图}{fig:cache-hardware-pie}

对\Case{080}，固定B计划的工期由113455降至90408 cycle，$S_{\rm hw}=1.2549$；C内重新选择后降至88140 cycle，$S_{\rm adapt}=1.0257$。对\Case{055}，固定计划出现189952字节命中，但B、C两侧工期均为25635 cycle，命中没有落在决定尾部完成时刻的关键路径上。图\ref{fig:cache-hit-gain}给出五核100图的$C(X_B^*)$命中率与同计划工期改善，零收益点保留在图中。
\samplefigure{cache_hit_gain_scatter}{五核字节命中率与同计划工期改善的逐图关系}{fig:cache-hit-gain}

\begin{table}[H]\centering\caption{四张代表图的五核三段工期与Cache访问（raw final\_selection）}\label{tab:c-cases}
\begin{tabular}{lrrrrr}\toprule
图 & $F_B(X_B^*)$ & $F_C(X_B^*)$ & $F_C(X_C^*)$ & $C(X_B^*)$命中/B & $S_{\rm hw}$\\\midrule
\Case{026} & 41665 & 41665 & 41665 & 95664 & 1.0000\\
\Case{053} & 492232 & 492232 & 491765 & 3570304 & 1.0000\\
\Case{055} & 25635 & 25635 & 25635 & 189952 & 1.0000\\
\Case{080} & 113455 & 90408 & 88140 & 2363392 & 1.2549\\\bottomrule
\end{tabular}\end{table}
表\ref{tab:c-cases}显示三种情形：\Case{055}读来源变化但完成时间不变；\Case{080}固定计划的硬件效应和C内适配均明显；\Case{053}命中字节较多但适配收益只有0.095\%，说明命中量与时间收益不成比例。

把同计划比较扩展到全部500个图--核组合，图\ref{fig:cache-hit-gain-3d}按核数分面，在二维坐标中表示$C(X_B^*)$字节命中率与$100[1-F_C(X_B^*)/F_B(X_B^*)]$工期降幅。326组最终记录包含命中，其中140组同计划工期未缩短；174组既无命中也无同计划工期变化。三条初始计划哈希不一致的记录保留在原始散点中，并不用于严格的同计划因果解释。
\samplefigure{cache_hit_gain_3d}{500组同计划配对的分核命中率与工期降幅分面散点图}{fig:cache-hit-gain-3d}

以哈希一致的五核98组为例，硬件效应改善、适配改善和综合改善分别覆盖48、37和64组。\Case{080}的C初始工期为90408 cycle，选中调度候选后降为88140 cycle；其余多数图的C内适配收益集中在零点附近。对\Case{040}，B最终工期为140837 cycle，C初始为140811 cycle，最终为140543 cycle，命中字节为12288；该记录属于500组原始配对中的同计划差异边界，三段工期分别反映B候选、C初始记录和C最终选择。

\begin{table}[H]\centering\small
\caption{五核哈希一致配对的硬件效应与调度适配交叉分布}
\label{tab:c-effect-cross}
\begin{tabular}{lrr}\toprule
效应分类 & 图数 & 综合工期降幅均值/\%\\\midrule
仅同计划硬件效应 & 27 & 2.8752\\
仅C内调度适配 & 16 & 3.9141\\
两种效应兼有 & 21 & 4.9809\\
两种效应均无 & 34 & 0.0000\\\bottomrule
\end{tabular}
\end{table}
严格同计划的98组中，硬件效应改善的48组与适配改善的37组有21组重合，综合改善为64组；其余34组的三段工期相同。该分解把Cache硬件、C内调度与计划来源差异分开记录。

\subsection{问题三结果}
问题三的每个图--核数组合均从该组合实际生成并通过官方评价的候选中按工期、额外搬运字典序选出$X_C^*$。相对整图单核基准，五核平均速度比为3.5533，中位数为4.2035；1500条主矩阵记录全部通过官方评价，五核结果中没有图慢于整图单核基准。逐核均值与首个成功候选到最终选择的平均改善见表\ref{tab:c-results}。
\begin{table}[H]\centering\caption{问题三：100图逐核速度比与候选搜索改善（raw final\_selection）}\label{tab:c-results}
\begin{tabular}{crrrr}\toprule
核数 & 平均速度比 & 中位数 & 慢于整图单核图数 & 首个成功至最终平均改善\\\midrule
1 & 1.0739 & 1.0073 & 0 & 1.0777\%\\
2 & 1.7998 & 1.9949 & 0 & 1.3230\%\\
3 & 2.4462 & 2.8901 & 0 & 1.1666\%\\
4 & 3.0448 & 3.6732 & 0 & 0.8289\%\\
5 & 3.5533 & 4.2035 & 0 & 1.1744\%\\\bottomrule
\end{tabular}\end{table}
\samplefigure{speedup_C}{问题三加速比均值、中位数及均值95\%重采样区间}{fig:speedup-c}

图\ref{fig:speedup-c}的浅色带为100张图有放回重采样1500次得到的均值95\%区间。raw final\_selection反映每个核数独立搜索的结果，因此相邻核数不强制单调：由4核增至5核时，72张图工期缩短、19张持平、9张变长，但全体平均速度比仍由3.0448升至3.5533。该差异来自不同核数下候选计划集合和有限评价预算的变化。

\subsection{问题三灵敏度分析：L2容量与带宽}
固定六张代表图的五核B计划，分别把L2 Cache容量与读带宽减半、加倍，比较固定计划和在扰动硬件下重新搜索的结果。固定计划实验共24组，全部通过官方Problem-3评估；重优化记录37条，37条均通过评估，每种变体均有3/6组改变计划。固定计划的工期比在四种扰动下均为1，重优化后的工期比与命中率见表\ref{tab:cache-sensitivity}。
\begin{table}[H]\centering\small\caption{六图四种Cache扰动的固定计划与重优化结果}\label{tab:cache-sensitivity}
\begin{tabular}{lrrrrr}\toprule
配置 & 图数 & 固定/默认工期比 & 重优化/默认工期比 & 固定命中率 & 重优化命中率\\\midrule
容量减半 & 6 & 1.0000 & 0.9900 & 2.10\% & 3.44\%\\
容量加倍 & 6 & 1.0000 & 0.9869 & 8.11\% & 13.02\%\\
带宽减半 & 6 & 1.0000 & 0.9892 & 4.78\% & 7.28\%\\
带宽加倍 & 6 & 1.0000 & 0.9885 & 4.78\% & 7.35\%\\\bottomrule
\end{tabular}\end{table}
表中的工期比先逐图计算，再在六张图间取等权平均；命中率同样按逐图字节命中率平均。重优化相对默认工期下降幅度依次为1.00\%、1.31\%、1.08\%和1.15\%，均来自\Case{016}、\Case{062}和\Case{067}三张图的计划调整。\Case{012}、\Case{051}和\Case{081}四种扰动下均保持原计划。

\begin{table}[H]\centering\small\caption{六张代表图重优化后的Cache参数逐例响应（相对默认工期比）}\label{tab:c-sensitivity-cases}
\begin{tabular}{lrrrr}\toprule
图 & 容量减半 & 容量加倍 & 带宽减半 & 带宽加倍\\\midrule
\Case{012} & 1.0000 & 1.0000 & 1.0000 & 1.0000\\
\Case{016} & 1.0000 & 1.0000 & 1.0000 & 1.0000\\
\Case{051} & 1.0000 & 1.0000 & 1.0000 & 1.0000\\
\Case{062} & 0.9758 & 0.9758 & 0.9758 & 0.9758\\
\Case{067} & 0.9643 & 0.9454 & 0.9594 & 0.9552\\
\Case{081} & 1.0000 & 1.0000 & 1.0000 & 1.0000\\\bottomrule
\end{tabular}\end{table}

容量和带宽的作用路径不同。\Case{051}容量加倍后固定计划命中率由8.33\%升至16.67\%，但重优化仍选回原计划，工期保持541017 cycle；\Case{062}在四种扰动下都由原计划切换为分割候选，重优化工期约为默认的0.9758，命中率随容量由1.74\%至6.16\%变化；\Case{067}固定计划命中为零，重优化引入分割后产生5.34\%至23.10\%的命中率，容量加倍对应工期比最低，为0.9454。命中率的变化通过查询服务池和计划结构共同作用于尾部完成时间。

\paragraph{小结。}
在六图四种扰动的固定计划实验中，工期保持不变；重新搜索后每种变体均有3张图改变计划，重优化/默认工期比介于0.9869和0.9900之间。L2容量和带宽对切图调度的影响集中在少数容量压力较高的图，必须通过计划重搜索才能显现。
